\scanpage{19}
Le $\unsure{\bar m|\widetilde m|\text{cas}}$ pb peut se formuler pour un
système de cycles

\[
\left\{\begin{aligned}
&Z_1^{\deleted{\unsure{+|*}}},\ldots,Z_k\qquad\text{de codim }c_1,\ldots c_k\\
&c_1+\ldots+c_k=n+1\\
&\text{\inserted{deux parmi les termes}}\ Z_{i\eta}\ \text{\inserted{sont} alg. équiv. à }0\\
&Z_1.\ldots.Z_k\quad\text{défini}
\end{aligned}\right.
\]

On veut pouvoir écrire

\[
 i(Z_1.\ldots.Z_k)=(Z_{1\eta},\ldots,Z_{k\eta})_v
\]

\uline{Propriétés du symbole local}

(i) Bilinéarité (\unsure{partielle|partielles} \ldots), symétrie

(ii) Invariance par isomorphismes sur $X_\eta$

(iibis) Pour $k=2$, formule de transposition
pour un morphisme $X_\eta\unsure{\text{---}|\longrightarrow}Y_\eta$

(iii) \deleted{\unsure{Bilin.|Bilinéarité}} \unsure{Normalisations|Normalisation} Si $Z_{k\eta}=Z'_{k\eta}\operatorname{div}f_\eta$,
\deleted{\unsure{avec|Ou}} où $Z_{1\eta}.\ldots Z_{k-1,\eta}.Z'_{k\eta}$ est défini (i.e.
l’intersection des $\operatorname{supp}Z_{1\eta}\ldots\operatorname{supp}Z_{k-1,\eta}$, $\operatorname{supp}Z'_{k\eta}$
est discrète] on a

\[
 \deleted{\unsure{i_v}}\ i_v(Z_{1\eta},\ldots,Z_{k\eta})=v(f_\eta(Z_{1\eta}.\ldots.Z_{k-1,\eta}.Z'_{k\eta} \unsure{)|,|))}
\]

(iv) Principe de continuité \ldots{} .

% margin note written sideways
\margin{(V) Extension du corps de base,}
% archive sheet number: 18
